\subsection{Fewer finite choices and smaller algebra rank}
\label{subsec:shared-degree-refinements}

The following two refinements preserve the retained-coordinate counts.
They improve the subsequent degree bound by removing unnecessary
finite choices and finite algebra generators.  They apply in
particular to the forty-three-coordinate construction.

\begin{lemma}[Fixed nonzero height scales]
\label{lem:fixed-height-scales}
Write $x(Q)=x_0+x_1\sqrt2$ for any nonidentity root point of the
shared construction.  One may replace the two piecewise height
scales attached to these components by
\[
                         A_0=x_0^2,
              \qquad A_1=(1+x_1^2)^2.
\]
Both are positive nonzero rational squares.  This removes all
eight zero-component case parameters and their closed equations,
without changing the arithmetic soundness and completeness proofs.
The original point-height bound with constant $7H$ still holds.
\end{lemma}
\begin{proof}
The point $Q$ is the double of a halfpoint $(a,b)$ on
$y^2=x^3-Dx$, with $D=(3-\sqrt2)^2$ and $b\ne0$.
The duplication formula gives
\[
                   x(Q)=\left(\frac{a^2+D}{2b}\right)^2.
\]
At each real embedding the numerator $a^2+D$ is strictly positive
and the denominator is nonzero.  Thus both real embeddings of
$x(Q)$ are positive, and their average $x_0$ is positive.
The assertions about the two scales follow.

Put $T=(H+\log3)|N|^\alpha$, where $\alpha$ is the radical
exponent used in the applicable height comparison.  The common
slope detector and each scaled detector give
\[
               h(\lambda),h(A_i\lambda)\le T,
               \qquad h(A_i)\le2T.
\]
For the first component $h(A_0)=2h(x_0)$, so $h(x_0)\le T$.
For a rational number $x_1=a/b$ in lowest terms with $b>0$,
\[
 h(1+x_1^2)=\log(a^2+b^2)\ge2h(x_1),
 \qquad h(A_1)\ge4h(x_1).
\]
The displayed equality uses
$\gcd(a^2+b^2,b^2)=1$.  Therefore $h(x_1)\le T/2$.  The
standard height inequality for the fixed quadratic basis gives
\[
 h(x(Q))\le h(x_0)+h(x_1)+\tfrac32\log2
 \le\tfrac32(H+\log3)|N|^\alpha+\tfrac32\log2
 <7H|N|^\alpha,
\]
because $H\ge1$ and $|N|\ge1$.  Thus the soundness estimates are
preserved, with room in the same constant.  Every completeness
argument uses only that the scales are fixed positive rational
squares once the root points are chosen; the new choices have
exactly those properties.  All nine detector constructions and
their simultaneous prime patterns remain available.
\end{proof}

\begin{lemma}[Parametrizing the first pole conic]
\label{lem:pole-conic-parametrization}
Let $B\in\mathbb Q^*$ and consider
\[
                        p^2=B(1-e)(1+3e).
\]
Its rational points are covered by two rational formulas with
one parameter.  For a finite choice $\sigma\in\{0,1\}$ retain
one rational coordinate $t$, put $D_t=t^2+3B$, require $D_t\ne0$,
and set
\[
 e=\frac{t^2+(4\sigma-1)B}{D_t},
 \qquad p=-\frac{4Bt(1-\sigma)}{D_t}.
\]
Every permitted assignment satisfies the conic equation, and every
old rational point occurs in one of the two choices.  In particular
the square root coordinate $p$ can be eliminated exactly.
\end{lemma}
\begin{proof}
For $\sigma=0$,
\[
                 1-e=\frac{4B}{D_t},\qquad
                 1+3e=\frac{4t^2}{D_t},
\]
which gives the equation.  For $\sigma=1$ the formulas give
$e=1,p=0$ directly.  Conversely, if an old point has $e\ne1$,
put $t=p/(e-1)$.  Substitution of its conic equation gives
\[
                       t^2+3B=-\frac{4B}{e-1}\ne0.
\]
The formulas with $\sigma=0$ recover both $e$ and $p$.  If
$e=1$, then $p=0$; the choice $\sigma=1,t=0$ gives this point
with denominator $3B\ne0$.  The other point with $p=0$,
namely $e=-1/3$, is already covered by $\sigma=0,t=0$.
Accordingly no nonzero guard on $p$ or on $t$ is imposed.
\end{proof}

Apply this lemma separately to the first conic in each of the
nine coordinate detectors and in the scalar block.  Their
$B$ values are nonzero by the existing guards.  Each block now
retains its conic parameter $t$ and its coordinate $s$, and
keeps only the second square root $q$ as a finite algebra
generator.  The parameter $t$ here parametrizes a conic; it
is distinct from the scalar preparation coordinate that was
removed in the preceding subsection.  The exceptional conic
choices are fixed finite parameters, not quantified selectors.

The preferred finite-anchor source therefore has $42$ retained
coordinates and twenty finite generators: nine quartic point
generators, ten quadratic pole generators, and one quadratic
Lucas generator.  Its algebra rank is at most
\[
                         4^9 2^{11}=2^{29}.
\]
The semilocal source has $44$ retained coordinates and one
additional quadratic anchor generator, hence rank at most
$2^{30}$.  Finite-algebra compression still gives totals
$43$ and $45$, respectively.  The ten conic choices replace
the ten now-absent first-root zero/nonzero strata, whereas
all eight zero-component choices have disappeared.  The
generic determinant compiler must use these updated ranks
and finite branch lists for its final degree bounds.
